% Tutorial, Chapter 7 -- Building a Forecast on Your Own: Lessons 25 to 28
% (AN-17).  A second worksheet, built from an empty worksheet with the
% skills of Chapters 2 to 6 and tied together: a growth rate referred to
% absolutely, copies of rows and columns, a range name made from a label,
% statistical @functions and @IF, and protection.  Figures: screens/tl25
% .. tl28b.k123 on the data volume of stages 4 to 7 (samples/5-fc25,
% 6-fc26, 7-fc27: the shipped application typed each lesson's keystrokes
% and saved FC25, FC26, FC27, the worksheet the next lesson retrieves).
% Sources (Law 1): os/i123/app/face_a.c (entries; /Copy with a one-cell
% TO; /Range Name Labels Right; /Range Name Table; /Range Format; /Range
% Unprot; /Worksheet Global Protection Enable and "Cell is protected";
% /File Save), os/i123/core/range.c (a name from a label; $NAME absolute),
% os/i123/core/eval.c (@SUM, @MAX, @MIN, @AVG, @IF), os/i123/core/fmt.c
% (the comma format; Percent).  Example: examples/forecast27, replayed.
% GATES: test-i123-formula test-i123-range-menu test-i123-copy-move-model test-flair-i123-files test-pubs-i123-examples test-pubs-screens test-pubs-i123
\Gates{test-i123-formula, test-i123-range-menu, test-i123-copy-move-model,
  test-flair-i123-files, test-pubs-i123-examples, test-pubs-screens,
  test-pubs-i123}
\chapter{Building a Forecast on Your Own}

In this chapter you build a second worksheet from the beginning: a sales
forecast for the four sales regions of Initech Systems Corporation.  The
first quarter's sales are known; the next three quarters are expected to
grow by a rate you set in one cell.  You use the skills of the earlier
chapters together, with fewer steps spelled out, and you see the reason
for building a worksheet this way: when you change the growth rate, the
whole forecast changes with it.

Start Initech 123, or erase the worksheet in memory with
\slashcmd{/Worksheet Erase Yes}.  Each lesson ends by saving the forecast,
and the next lesson begins by retrieving it, so you can stop between
lessons.

\FlairShot[fig:tu7open]{tl27}{The forecast as Lesson~27 leaves it}

\section{Lesson 25: Planning the Forecast}
\ix{planning a worksheet}
In this lesson you learn to:
\begin{itemize}
\item plan a worksheet so that its assumptions are in cells of their own
\item set the width of every column at once
\item enter a percentage
\end{itemize}

\subsection{The Plan}
A worksheet is easiest to change when every number it assumes is typed in
one place, and every other number is calculated from it.  The forecast has
three blocks:
\begin{LotTable}{p{30mm}p{24mm}p{70mm}}
\LotHead{\textbf{Item} & \textbf{Address} & \textbf{Meaning}}
Growth & B3 & The growth rate from one quarter to the next\\
Sales & B6..E9 & The sales of each region in each quarter; only the
  first quarter, B6..B9, is typed\\
Totals & F6..F10, B10..E10 & The total of each region for the year and
  of each quarter for all regions\LotEnd
\end{LotTable}

\subsection{The Labels and the Growth Rate}
The numbers of the forecast need columns a little wider than the standard
9 characters.

\begin{Keystrokes}
\Do{Press}{\kname{SLASH}}{and type \type{w}, \type{g}, \type{c} to select
  Worksheet Global Column-Width}
\Do{Type}{\type{10}}{and press \kname{ENTER}}
\Do{Type}{\type{Sales Forecast 1993}}{in A1, and press \kname{ENTER}}
\Do{Press}{\fkey{Goto}{F5}}{type \type{a3} and press \kname{ENTER}}
\Do{Type}{\type{Growth}}{and press \kname{RIGHT}}
\Do{Type}{\type{5\%}}{the growth rate, and press \kname{ENTER}}
\end{Keystrokes}
Initech 123 stores \type{5\%} as the number 0.05.  Now type the headings
and the regions:

\begin{Keystrokes}
\Do{Press}{\fkey{Goto}{F5}}{type \type{a5} and press \kname{ENTER}}
\Do{Type}{\type{Region}}{and press \kname{RIGHT}}
\Do{Type}{\type{"Q1}}{and press \kname{RIGHT}; then \type{"Q2},
  \kname{RIGHT}, \type{"Q3}, \kname{RIGHT}, \type{"Q4}, \kname{RIGHT}}
\Do{Type}{\type{"Year}}{and press \kname{ENTER}}
\Do{Press}{\fkey{Goto}{F5}}{type \type{a6} and press \kname{ENTER}}
\Do{Type}{\type{North}}{and press \kname{DOWN}; then \type{South},
  \kname{DOWN}, \type{East}, \kname{DOWN}, \type{West}, \kname{DOWN}}
\Do{Type}{\type{Total}}{and press \kname{ENTER}}
\end{Keystrokes}

\subsection{Formatting the Rate and Saving}
B3 shows 0.05.  Give it the Percent format with no decimal places, then
save the worksheet as FC25.

\begin{Keystrokes}
\Do{Press}{\kname{SLASH}}{and type \type{r}, \type{f}, \type{p}}
\Do{Type}{\type{0}}{no decimal places, and press \kname{ENTER}}
\Do{Type}{\type{b3}}{and press \kname{ENTER}}
\Do{Press}{\kname{SLASH}}{and type \type{f}, \type{s}, type \type{FC25}
  and press \kname{ENTER}}
\end{Keystrokes}
\FlairShot[fig:tl25]{tl25}{The plan of the forecast}

\section{Lesson 26: Growing the Numbers}
\ix{growth formula}
In this lesson you learn to:
\begin{itemize}
\item write a formula that refers to an assumption absolutely
\item copy a row of formulas down several rows at once
\item total rows and columns with copies of one formula
\end{itemize}

\subsection{The First Quarter}
Retrieve FC25, then type the first quarter's sales.

\begin{Keystrokes}
\Do{Press}{\kname{SLASH}}{type \type{f}, \type{r}, type \type{FC25} and
  press \kname{ENTER}}
\Do{Press}{\fkey{Goto}{F5}}{type \type{b6} and press \kname{ENTER}}
\Do{Type}{\type{1200}}{and press \kname{DOWN}; then \type{950},
  \kname{DOWN}, \type{1430}, \kname{DOWN}}
\Do{Type}{\type{780}}{and press \kname{ENTER}}
\end{Keystrokes}

\subsection{The Growth Formula}
Each quarter's sales are the quarter before, multiplied by one plus the
growth rate.  For North in the second quarter that is
\texttt{+B6*(1+\$B\$3)}: B6 is relative, so that each copy uses the
quarter to its left, and \texttt{\$B\$3} is absolute, so that every copy
uses the one growth rate.

\begin{Keystrokes}
\Do{Press}{\fkey{Goto}{F5}}{type \type{c6} and press \kname{ENTER}}
\Do{Type}{\type{+b6*(1+\$b\$3)}}{and press \kname{ENTER}}
\Do{Press}{\kname{SLASH}}{type \type{c}, press \kname{ENTER}, type
  \type{d6} and press \kname{ENTER}}
\Do{Press}{\kname{SLASH}}{type \type{c}, press \kname{ENTER}, type
  \type{e6} and press \kname{ENTER}}
\end{Keystrokes}
Row 6 now holds the forecast for North (Figure~\ref{fig:tl26a}).  The
control panel shows the formula of C6.
\FlairShot[fig:tl26a]{tl26a}{The growth formula copied across row 6}

\subsection{Copying a Row Down}
The three formulas of C6..E6 serve every region.  Copy them as a block:
when you give the upper left cell of the TO range, the whole row arrives.

\begin{Keystrokes}
\Do{Press}{\kname{SLASH}}{type \type{c}, type \type{c6..e6} and press
  \kname{ENTER}}
\Do{Type}{\type{c7}}{and press \kname{ENTER}}
\end{Keystrokes}
Copy C6..E6 to C8 and to C9 in the same way.

\subsection{The Totals}
\begin{Keystrokes}
\Do{Press}{\fkey{Goto}{F5}}{type \type{f6} and press \kname{ENTER}}
\Do{Type}{\type{@sum(b6..e6)}}{the year for North, and press
  \kname{ENTER}}
\Do{Press}{\kname{SLASH}}{type \type{c}, press \kname{ENTER}, type
  \type{f7} and press \kname{ENTER}}
\Do{Press}{\kname{SLASH}}{type \type{c}, type \type{f6..f7} and press
  \kname{ENTER}, type \type{f8} and press \kname{ENTER}}
\Do{Press}{\kname{SLASH}}{type \type{c}, type \type{f8} and press
  \kname{ENTER}, type \type{f9} and press \kname{ENTER}}
\Do{Press}{\fkey{Goto}{F5}}{type \type{b10} and press \kname{ENTER}}
\Do{Type}{\type{@sum(b6..b9)}}{the first quarter for all regions, and
  press \kname{ENTER}}
\end{Keystrokes}
Copy B10 to C10; copy B10..C10 to D10; copy B10..C10 to E10; and copy B10
to F10, where the formula adds the four yearly totals above it.

\subsection{Formatting and Saving}
\begin{Keystrokes}
\Do{Press}{\kname{SLASH}}{and type \type{r}, \type{f}, \type{,} (Range
  Format with commas)}
\Do{Type}{\type{0}}{and press \kname{ENTER}}
\Do{Type}{\type{b6..f10}}{and press \kname{ENTER}}
\Do{Press}{\kname{SLASH}}{and type \type{f}, \type{s}, type \type{FC26}
  and press \kname{ENTER}}
\end{Keystrokes}
The forecast is complete (Figure~\ref{fig:tl26b}).  The numbers are
displayed without decimal places, but each is calculated in full: the
totals add the full values, not the displayed ones.
\FlairShot[fig:tl26b]{tl26b}{The finished forecast at 5\% growth}

\section{Lesson 27: Naming the Rate and Summarizing}
\ix{range names}\ix{"@MAX}\ix{"@MIN}\ix{"@AVG}
In this lesson you learn to:
\begin{itemize}
\item name a cell after the label beside it
\item use the name in a formula
\item summarize a row with the statistical @functions
\item test the result against a target with \fn{@IF}
\end{itemize}

\subsection{Naming the Growth Rate}
The label \texttt{Growth} in A3 is just the name the growth rate needs.
\slashcmd{/Range Name Labels Right} names the cell to the right of each
label in a range.

\begin{Keystrokes}
\Do{Press}{\kname{SLASH}}{type \type{f}, \type{r}, type \type{FC26} and
  press \kname{ENTER}}
\Do{Press}{\kname{SLASH}}{and type \type{r}, \type{n}, \type{l}, \type{r}
  to select Range Name Labels Right}
\Do{Type}{\type{a3}}{the label, and press \kname{ENTER}}
\end{Keystrokes}
B3 is now named GROWTH.  Type the growth formula again with the name, and
copy it as before:

\begin{Keystrokes}
\Do{Press}{\fkey{Goto}{F5}}{type \type{c6} and press \kname{ENTER}}
\Do{Type}{\type{+b6*(1+\$growth)}}{and press \kname{ENTER}}
\end{Keystrokes}
Copy C6 to D6 and to E6, and C6..E6 to C7, C8 and C9, exactly as in
Lesson~26.  The control panel still shows \texttt{+B6*(1+\$B\$3)}: a name
stands for its range, and the dollar sign before \texttt{GROWTH} makes the
reference absolute.

\subsection{The Summary}
\begin{Keystrokes}
\Do{Press}{\fkey{Goto}{F5}}{type \type{a12} and press \kname{ENTER}}
\Do{Type}{\type{Best}}{and press \kname{DOWN}; then \type{Lowest},
  \kname{DOWN}, \type{Average}, \kname{DOWN}}
\Do{Type}{\type{Target?}}{and press \kname{ENTER}}
\Do{Press}{\fkey{Goto}{F5}}{type \type{b12} and press \kname{ENTER}}
\Do{Type}{\type{@max(b10..e10)}}{the best quarter, and press
  \kname{DOWN}}
\Do{Type}{\type{@min(b10..e10)}}{the lowest quarter, and press
  \kname{DOWN}}
\Do{Type}{\type{@avg(b10..e10)}}{the average quarter, and press
  \kname{DOWN}}
\Do{Type}{\type{@if(f10>=20000,"Met","Not met")}}{and press
  \kname{ENTER}}
\end{Keystrokes}
The target of the sales office is 20,000 for the year.  \fn{@IF} shows
\texttt{Met} when the total in F10 reaches it, and \texttt{Not met}
otherwise.

\subsection{A Table of the Names}
Format B12..B14 with commas and no decimal places, as in Lesson~26.  Then
list the names of the worksheet:

\begin{Keystrokes}
\Do{Press}{\fkey{Goto}{F5}}{type \type{a17} and press \kname{ENTER}}
\Do{Press}{\kname{SLASH}}{and type \type{r}, \type{n}, \type{t} to select
  Range Name Table}
\Do{Press}{\kname{ENTER}}{to place the table at A17}
\Do{Press}{\kname{SLASH}}{and type \type{f}, \type{s}, type \type{FC27}
  and press \kname{ENTER}}
\end{Keystrokes}
\FlairShot[fig:tl27]{tl27}{The summary, the target test and the table of
names}

At 5\% growth the year falls short of the target.  The formulas of the
forecast, and what they show at 5\% and at 10\% growth:

\LotExample{examples/forecast27}

\section{Lesson 28: Trying Another Rate and Protecting the Forecast}
\ix{what-if}\ix{protection}
In this lesson you learn to:
\begin{itemize}
\item change an assumption and watch the worksheet follow it
\item protect a worksheet so that only its assumptions can be changed
\end{itemize}

\subsection{Another Growth Rate}
\begin{Keystrokes}
\Do{Press}{\kname{SLASH}}{type \type{f}, \type{r}, type \type{FC27} and
  press \kname{ENTER}}
\Do{Press}{\fkey{Goto}{F5}}{type \type{b3} and press \kname{ENTER}}
\Do{Type}{\type{10\%}}{and press \kname{ENTER}}
\end{Keystrokes}
Every forecast cell, every total and the summary change at once.  At 10\%
growth the year reaches the target, and B15 shows \texttt{Met}.

\subsection{Protecting the Forecast}
Only the growth rate and the first quarter are typed; everything else is
calculated.  Unprotect the typed cells and enable protection, so that no
formula can be typed over by mistake.

\begin{Keystrokes}
\Do{Press}{\kname{SLASH}}{and type \type{r}, \type{u}; press
  \kname{ENTER} to unprotect B3}
\Do{Press}{\kname{SLASH}}{and type \type{r}, \type{u}; type \type{b6..b9}
  and press \kname{ENTER}}
\Do{Press}{\kname{SLASH}}{and type \type{w}, \type{g}, \type{p}, \type{e}
  to enable protection}
\Do{Press}{\fkey{Goto}{F5}}{type \type{c6} and press \kname{ENTER}}
\Do{Type}{\type{5}}{and press \kname{ENTER}}
\end{Keystrokes}
Initech 123 refuses the entry, because C6 holds a formula and is protected
(Figure~\ref{fig:tl28a}).
\FlairShot[fig:tl28a]{tl28a}{A forecast formula is protected}
\begin{Keystrokes}
\Do{Press}{\kname{ESC}}{to return to the worksheet}
\Do{Press}{\kname{SLASH}}{and type \type{f}, \type{s}, type
  \type{FORECAST} and press \kname{ENTER}}
\end{Keystrokes}
\FlairShot[fig:tl28b]{tl28b}{The protected forecast saved as FORECAST}

You can now try any growth rate, or new first-quarter sales, by typing
over B3 or B6..B9; nothing else can be changed until you disable
protection with \slashcmd{/Worksheet Global Protection Disable}.

This completes the \doctitle{Tutorial}.  The \doctitle{Initech 123
Reference} describes every command you have used, and all the others.
